% Authoritative LaTeX manuscript; Markdown is a generated compatibility projection.

MAGE aims to increase the reliable engineering work accomplished with a given investment of human effort and supporting machinery. The cost-per-correct-admitted-change equation in \autoref{6.1-toward-a-theory-of-mage:heading:cost-and-failure-separately} develops this economic account; \autoref{chapter:5.3-evaluating-docable} examines DocAble's costs and outcomes. The practical question is whether introducing models and controls improves those outcomes enough to justify construction and maintenance.

\subsection{Metrics for Engineering Performance}\label{heading:metrics-for-engineering-performance}

\autoref{practicing-primary-readings} identifies five useful metrics. They measure different properties and should not be combined into a maturity score. Higher throughput does not compensate for unacceptable defects, and additional model coverage is useful only when it supports consequential engineering questions. Model quality and return on engineering capital may require several indicators or case-based assessments; report their basis rather than implying a single numerical measure.

{\footnotesize
\begin{longtable}[]{@{}p{.22\linewidth}p{.35\linewidth}p{.33\linewidth}@{}}
\caption[Metrics for evaluating MAGE]{\emph{Metrics for evaluating MAGE.} The metrics connect engineering outcomes to the book's account of Modeling, Alignment, and engineering capital. They measure different properties and should be interpreted together rather than combined into a single score.}\label{practicing-primary-readings}\\
\toprule
\textbf{Metric} & \textbf{Measurement} & \textbf{Engineering significance} \\
\midrule
\endhead
\textbf{Durable throughput} & Accepted work units per unit time, distinguishing work that survives a defined observation period from work requiring subsequent repair or rollback. & Measures cumulative progress rather than implementation activity. \\
\textbf{Defect escape rate} & Violations discovered after the enforcement boundary intended to prevent them, relative to relevant accepted work. & Evaluates whether Alignment mechanisms prevent unacceptable work from proceeding. \\
\textbf{Human engineering effort} & Human time spent specifying, reviewing, investigating, correcting, and supervising work, relative to durable output. & Measures how much scarce human effort delegated implementation continues to require. \\
\textbf{Model quality and coverage} & Coverage of selected engineering questions and obligations; correspondence failures, unresolved drift, and validation results. & Evaluates whether models provide trustworthy representations of consequential system properties. \\
\textbf{Return on engineering capital} & Subsequent engineering effort or losses avoided through reusable models and controls, compared with their construction and maintenance costs. & Evaluates whether governance conversion produces lasting economic value. \\
\bottomrule
\end{longtable}
}

\textbf{Establish a baseline.} Before changing the environment, record comparable work units, their acceptance criteria, completion evidence, human interventions, and subsequent repairs. Use the same unit and observation window afterward. Report differences in complexity or consequence rather than attributing every change to MAGE. A before-and-after comparison can suggest whether the investment helped; it does not by itself establish causation.

Agree on the unit, boundary, and observation window before interpreting a trend. A landed change, completed work order, and released capability are different units. Keep their denominators visible and distinguish changes in the workload from changes in the machinery. Retain the revisions and evidence behind consequential observations. A dashboard can help locate a change in behavior; it cannot supply a comparison that the underlying records do not support.

\subsection{Evaluating Outcomes and Costs}\label{heading:evaluating-outcomes-and-costs}

\textbf{Durable throughput.} Track reopenings, rollbacks, and repeated repairs within the observation window. Report recently accepted work whose durability is not yet observable separately from work that has survived the window. A later requirements change does not necessarily invalidate the original result; repeated repair of the original obligation does. Commit counts and immediate ticket closure omit this distinction. This measures the cumulative progress distinguished from implementation activity in \autoref{6.1-toward-a-theory-of-mage:heading:the-dynamic-model}.

\textbf{Defect escape rate.} Record the violated obligation, intended enforcement boundary, affected revision, and discovery. A failed test before admission demonstrates detection; the same defect found after that boundary demonstrates an escape. For a work-unit escape rate, divide accepted units with an escape by relevant accepted units in the same cohort and observation window; report defect counts separately. Compare rates only when obligations and opportunities for observation are sufficiently comparable. More reports may reflect better observation, and no reports may reflect weak observation. The acceptance mechanisms in \autoref{chapter:3.3-constraints-sensors-validators-gates} supply the boundaries whose effectiveness is being evaluated.

\textbf{Human engineering effort.} Record specification, repeated reconstruction, review, exception handling, recovery, and consequential judgment, using time where practical and intervention records otherwise. Relate effort to durable output. Keep intentional engineering decisions visible: choosing an obligation or interpreting ambiguous evidence differs from explaining the same convention again. Effort has been saved when work becomes unnecessary or the environment performs it reliably, not merely when it stops being recorded.

\textbf{Return on engineering capital.} Retain concrete examples of later benefit: a dependency no longer reconstructed, a recurring failure prevented, a larger task made tractable, or knowledge retained after its original author has left. Account for construction, execution, understanding, reconciliation, and maintenance over a declared period. A model may repay its cost without preventing a defect; a control may justify itself through a consequential risk rather than frequent use. Reuse counts alone establish neither return. A temporary increase in effort or decrease in throughput may fund later capability; evaluate that investment across subsequent work, not only its first use.

\subsection{Evaluating the Engineering Environment}\label{heading:evaluating-the-engineering-environment}

\textbf{Model quality and coverage} depend on the claims each model makes. A descriptive model must correspond to its subject; a normative model states intent the system must satisfy; a derived view depends on its source and transformation. Historical work records have the different lifetime explained in \href{appendix-b-getting-started.html\#use-a-change-as-the-first-model}{``Model the Change'' in B.1}. Examine drift findings, failed queries, missing consequential coverage, and reconciliation evidence. Freshness means correspondence after relevant change, not a recent edit date or valid link. \autoref{2.6-system-knowledge:heading:maintaining-correspondence} develops these maintenance obligations.

Correspondence alone does not establish that the model expresses the right requirement, as \autoref{3.2-models-and-the-semantic-gap:heading:agreement-is-not-correctness} explains. Check the machinery as well: a validator may inspect the wrong artifact, or a gate may fail to act on a finding. These failures undermine both acceptance and the metrics used to evaluate it.

Include the health of observation in the review. Did scheduled checks run? Which revisions did they inspect? Could a check fail without anyone seeing or acting on its result? Were exceptions closed through verified repair or only acknowledgment? An empty findings list is uninterpretable when the check never ran. Similarly, an uncalibrated visual indicator is status, not a quantitative measurement. Preserve enough of the observation process to distinguish a quiet system from a silent sensor.

\subsection{Diagnosing Changes in Performance}\label{heading:diagnosing-changes-in-performance}

Review at a cadence suited to the work and after consequential incidents. Use diagnostic metrics to investigate changes in outcomes and costs. When throughput declines, inspect queue length, validation latency, and human intervention. When escapes increase, examine model coverage, checker coverage, and enforcement failures. When human effort rises, distinguish repeated reconstruction, review, and recovery. Churn may reveal repeated repair or useful restructuring; closure rate may conceal reopened work. A support ratio describes surrounding machinery relative to implementation, not its value. More controls may indicate protection or interference. Investigate deteriorating model quality through drift, broken traceability, and missing correspondence checks. Falling capital return calls for examining carrying costs and obsolete controls. The observations in \autoref{chapter:5.3-evaluating-docable} are evidence from particular settings, not universal operating targets.

Use these measurements to identify where engineering effort is being lost and whether existing controls are effective. Repeated reconstruction may justify a shared model; recurring defects may justify stronger validation or enforcement; excessive validation cost may justify redesigning the checking machinery. These are candidates for governance conversion, as developed in \autoref{chapter:3.4-governance-conversion}, not automatic prescriptions. After making a change, measure again to determine whether the new engineering structure improved subsequent work. Retire models and controls whose costs exceed their continuing value.
